\chapter*{Simulator Line Commands}
\addcontentsline{toc}{chapter}{Simulator Line Commands}

This chapter presents the line mode commands. Although they are called line mode commands, these commands are also available when windows are enabled. The following table is summary of the available commands. Like all commands, they are entered in the command window.

\section*{Command Summary}
\addcontentsline{toc}{section}{Command Summary}

\begin{longtable}{@{}p{0.3\linewidth}p{0.6\linewidth}@{}}
    \caption{Simulator Line Commands} \\
    \toprule
    \textbf{Command} & \textbf{Purpose} \\
    \midrule
    \endfirsthead
    \toprule
     \textbf{Command} & \textbf{Purpose} \\
    \midrule
    \endhead
    \midrule
    \multicolumn{2}{r}{\textit{Continued on next page}} \\
    \midrule
    \endfoot
    \bottomrule
    \endlastfoot
    \textbf{DA} & display absolute memory \\
	\midrule
	\textbf{DM} & display configured Twin-64 modules \\
	\midrule
	\textbf{DO} & execute a command from the command history stack \\
	\midrule
	\textbf{DW} & display windows \\
	\midrule
	\textbf{ENV} & display, add, modify environment variables \\
	\midrule
	\textbf{EXIT} & leave the Twin-64 simulator \\
	\midrule
	\textbf{HELP} & display help information \\
	\midrule
	\textbf{HIST} & display command history stack \\
	\midrule
	\textbf{ITLB} & insert into the unified TLB of a CPU \\
	\midrule
	\textbf{LF} & load ELF file into memory \\
	\midrule
	\textbf{MA} & modify absolute memory \\
	\midrule
	\textbf{MR} & modify a register of a CPU \\
	\midrule
	\textbf{NM} & add a module to the simulated environment \\
	\midrule
	\textbf{PTLB} & purge from the unified TLB of a CPU \\
	\midrule
	\textbf{REDO} & edit and do a command from the command history stack \\
	\midrule
	\textbf{RM} & remove a module from the simulated environment \\
	\midrule
	\textbf{RESET} & reset a Twin64 system or CPU \\
	\midrule
	\textbf{RUN} & run a configured simulation \\
	\midrule
	\textbf{STEP} & step through a simulation\\
	\midrule
	\textbf{W} & evaluate and display an expression \\
	\midrule
	\textbf{XF} & execute commands from a file \\
	\end{longtable}
	

\section*{Command Syntax}
\addcontentsline{toc}{section}{Command Syntax}

Command syntax is described using the following conventions. Reserved words
are written in uppercase letters. Parameters are written in lowercase letters.

Optional elements are enclosed in square brackets \texttt{[ ]}. Elements that
may be repeated are enclosed in curly braces \texttt{\{ \}}.

Numeric values are written in decimal notation by default. When expressed in
hexadecimal notation, numeric values are prefixed with \texttt{0x}. For improved
readability, underscore characters may be used as digit separators within
hexadecimal values.

String literals are enclosed in double quotation marks (\texttt{" "}). Longer
strings may be constructed by placing multiple string literals adjacent to
each other.

A command line may extend across multiple lines. To continue a command on the
next line, place a backslash (\texttt{\textbackslash}) at the end of the line
before the carriage return.

The following sections describe each command in detail, including its purpose,
syntax, parameters, and usage examples.


%--------------------------------------------------------------------------------
%
%
%--------------------------------------------------------------------------------
\pagebreak
\section*{DA - display physical memory}
\addcontentsline{toc}{section}{DA - display physical memory}

%--------------------------------------------------------------------------------
\begin{iPageDescEntry}{Syntax}
\texttt{DA } \\
\end{iPageDescEntry}

%--------------------------------------------------------------------------------
\begin{iPageDescEntry}{Description}

The \texttt{DA} command ...

\end{iPageDescEntry}

%--------------------------------------------------------------------------------
\begin{iPageDescOpEntry}{Example}
\begin{lstlisting}[style=iPageOpStyle]

... 

\end{lstlisting}
\end{iPageDescOpEntry}

%--------------------------------------------------------------------------------
\begin{iPageDescEntry}{Notes}
None.
\end{iPageDescEntry}

%--------------------------------------------------------------------------------
%
%
%--------------------------------------------------------------------------------
\pagebreak
\section*{DM - list configured modules}
\addcontentsline{toc}{section}{DM - list configured modules}

%--------------------------------------------------------------------------------
\begin{iPageDescEntry}{Syntax}
\texttt{DM } \\
\end{iPageDescEntry}

%--------------------------------------------------------------------------------
\begin{iPageDescEntry}{Description}

The \texttt{DM} command ...

\end{iPageDescEntry}

%--------------------------------------------------------------------------------
\begin{iPageDescOpEntry}{Example}
\begin{lstlisting}[style=iPageOpStyle]

... 

\end{lstlisting}
\end{iPageDescOpEntry}

%--------------------------------------------------------------------------------
\begin{iPageDescEntry}{Notes}
None.
\end{iPageDescEntry}

%--------------------------------------------------------------------------------
%
%
%--------------------------------------------------------------------------------
\pagebreak
\section*{DO - perform command from history stack}
\addcontentsline{toc}{section}{DO - perform command from history stack}

%--------------------------------------------------------------------------------
\begin{iPageDescEntry}{Syntax}
\texttt{DO } \\
\end{iPageDescEntry}

%--------------------------------------------------------------------------------
\begin{iPageDescEntry}{Description}

The \texttt{DO} command ...

\end{iPageDescEntry}

%--------------------------------------------------------------------------------
\begin{iPageDescOpEntry}{Example}
\begin{lstlisting}[style=iPageOpStyle]

... 

\end{lstlisting}
\end{iPageDescOpEntry}

%--------------------------------------------------------------------------------
\begin{iPageDescEntry}{Notes}
None.
\end{iPageDescEntry}

%--------------------------------------------------------------------------------
%
%
%--------------------------------------------------------------------------------
\pagebreak
\section*{DW — List Window Info}
\addcontentsline{toc}{section}{DW — List Window Info}

%--------------------------------------------------------------------------------
\begin{iPageDescEntry}{Syntax}
\texttt{dw [ <wNum> ]}
\end{iPageDescEntry}

%--------------------------------------------------------------------------------
\begin{iPageDescEntry}{Description}

The \texttt{DW} command displays a list of all currently created windows.
Each window is shown with its name and type, window ID, assigned stack number,
and—if applicable—the associated module number and module type.

Windows are listed regardless of whether they are enabled or disabled.
If a window number is specified, only the selected window is displayed.

Specifying an optional stack number restricts the output to windows belonging
to that stack.

\end{iPageDescEntry}

%--------------------------------------------------------------------------------
\begin{iPageDescOpEntry}{Example}

The following example shows windows created using the \texttt{WN} command.
In this example, the module number for the CPU and TLB is \texttt{3}, and the
memory window starts at address \texttt{0x1000}.

\begin{lstlisting}[style=iPageOpStyle]
(5)  -> wn cpu, 3
(6)  -> wn ITLB, 3
(7)  -> wn MEM, 0x1000
(8)  -> dw
Name      Stack     Id        WType     Mod       MType     
CPU       1         1         CPU       3         PROC      
I-TLB     1         2         TLB       3         PROC      
MEM       2         3         Memory    1         MEM       
(9)  ->
(9)  -> dw 2
MEM       2         3         Memory    1         MEM 
(10) ->
\end{lstlisting}

\end{iPageDescOpEntry}

%--------------------------------------------------------------------------------
\begin{iPageDescEntry}{Notes}
See also the \texttt{WN} command.
\end{iPageDescEntry}

%--------------------------------------------------------------------------------
%
%
%--------------------------------------------------------------------------------
\pagebreak
\section*{ENV - manage environment variables}
\addcontentsline{toc}{section}{ENV - manage environment variables}

%--------------------------------------------------------------------------------
\begin{iPageDescEntry}{Syntax}
\texttt{ENV } \\
\end{iPageDescEntry}

%--------------------------------------------------------------------------------
\begin{iPageDescEntry}{Description}

The \texttt{ENV} command ...

\end{iPageDescEntry}

%--------------------------------------------------------------------------------
\begin{iPageDescOpEntry}{Example}
\begin{lstlisting}[style=iPageOpStyle]

... 

\end{lstlisting}
\end{iPageDescOpEntry}

%--------------------------------------------------------------------------------
\begin{iPageDescEntry}{Notes}
None.
\end{iPageDescEntry}

%--------------------------------------------------------------------------------
%
%
%--------------------------------------------------------------------------------
\pagebreak
\section*{EXIT - leave simulator}
\addcontentsline{toc}{section}{EEXIT - leave simulator}

%--------------------------------------------------------------------------------
\begin{iPageDescEntry}{Syntax}
\texttt{EXIT } \\
\end{iPageDescEntry}

%--------------------------------------------------------------------------------
\begin{iPageDescEntry}{Description}

The \texttt{EXIT} command ...

\end{iPageDescEntry}

%--------------------------------------------------------------------------------
\begin{iPageDescOpEntry}{Example}
\begin{lstlisting}[style=iPageOpStyle]

... 

\end{lstlisting}
\end{iPageDescOpEntry}

%--------------------------------------------------------------------------------
\begin{iPageDescEntry}{Notes}
None.
\end{iPageDescEntry}

%--------------------------------------------------------------------------------
%
%
%--------------------------------------------------------------------------------
\pagebreak
\section*{HELP - list help info}
\addcontentsline{toc}{section}{HELP - list help info}

%--------------------------------------------------------------------------------
\begin{iPageDescEntry}{Syntax}
\texttt{HELP } \\
\texttt{HELP <cmdId> } \\
\texttt{HELP "commands" } \\
\texttt{HELP "wcommands" } \\ 
\texttt{HELP "predefined" }
\end{iPageDescEntry}

%--------------------------------------------------------------------------------
\begin{iPageDescEntry}{Description}

The \texttt{HELP} command displays help information for the simulator. Help topics are grouped into several categories:

\begin{itemize}
  \item The \texttt{COMMANDS} keyword lists a summary of all available line-mode commands. Entering \texttt{HELP} with no arguments also displays this list.
  \item The \texttt{WCOMMANDS} keyword lists a summary of all available window-mode commands.
  \item The \texttt{PREDEFINED} keyword lists all predefined functions that can be used in command-line expressions.
\end{itemize}

To view detailed syntax for a specific command or predefined function, type \texttt{HELP <cmdId>}
where \textit{cmdId} is the name of the command or function.


\end{iPageDescEntry}

%--------------------------------------------------------------------------------
\begin{iPageDescOpEntry}{Example}
\begin{lstlisting}[style=iPageOpStyle]
(8) ->help
 help            list help info
 exit            program exit
 hist            command history
 do              re-execute command
 redo            edit and then re-execute command
 env             lists the env tab, a variable, sets a variable
 xf              execute commands from a file
 reset           resets the system, etc.
 run             run the system ( all CPUs )
 step            single step the system
 w               evaluates and prints an expression
 dm              display registered modules for the system
 mr              modify a CPU register
 da              display absolute memory
 ma              modify absolute memory
 pica            purges instruction cache line data
 pdca            purges data cache line data
 fdca            flushes data cache line data
 iitlb           inserts an entry into the instruction TLB
 idtlb           inserts an entry into the data TLB
 pitlb           purges an entry from the instruction TLB
 pdtlb           purges an entry from the data TLB
 won             enables windows mode / redraw
(9) ->
(9) ->help da
da <adr> [ , <len> ] [ , <fmt> ] - display absolute memory
(10) ->
\end{lstlisting}
\end{iPageDescOpEntry}

%--------------------------------------------------------------------------------
\begin{iPageDescEntry}{Notes}
None.
\end{iPageDescEntry}

%--------------------------------------------------------------------------------
%
%
%--------------------------------------------------------------------------------
\pagebreak
\section*{HIST - list command history}
\addcontentsline{toc}{section}{HIST - list command history}

%--------------------------------------------------------------------------------
\begin{iPageDescEntry}{Syntax}
\texttt{HIST } \\
\end{iPageDescEntry}

%--------------------------------------------------------------------------------
\begin{iPageDescEntry}{Description}

The \texttt{HIST} command ...

\end{iPageDescEntry}

%--------------------------------------------------------------------------------
\begin{iPageDescOpEntry}{Example}
\begin{lstlisting}[style=iPageOpStyle]

... 

\end{lstlisting}
\end{iPageDescOpEntry}

%--------------------------------------------------------------------------------
\begin{iPageDescEntry}{Notes}
None.
\end{iPageDescEntry}

%--------------------------------------------------------------------------------
%
%
%--------------------------------------------------------------------------------
\pagebreak
\section*{ITLB — Insert TLB Entry}
\addcontentsline{toc}{section}{ITLB — Insert TLB Entry}

%--------------------------------------------------------------------------------
\begin{iPageDescEntry}{Syntax}
\texttt{itlb <vAdr> , <pAdr> , <size> , <acc> [ , L [ , U ]]}
\end{iPageDescEntry}

%--------------------------------------------------------------------------------
\begin{iPageDescEntry}{Description}

The \texttt{ITLB} command inserts a translation entry into the unified Translation Lookaside Buffer (TLB). The \texttt{vAdr} parameter specifies the virtual address of the mapping, while \texttt{pAdr} specifies the corresponding physical memory address.

To reduce TLB lookup overhead and improve translation efficiency, the TLB supports mappings over contiguous virtual address ranges rather than only fixed-size pages. Multiple page sizes are supported, allowing software to choose an appropriate granularity based on usage patterns. The supported page sizes are 4~KB, 64~KB, 1~MB, and 16~MB.

The \texttt{acc} parameter encodes the access permissions and privilege requirements for the translation as shown in the TLB description. It specifies whether the mapped page may be read, written, or executed, as well as whether it is a gateway page along with the minimum privilege level required for data fetch. Every data fetch is checked against these access attributes, and violations result in a protection fault trap.

Optional flags may be supplied to control the behavior of the TLB entry. These flags determine how the entry participates in replacement and caching decisions.

\begin{smallGapItemize}
  \item \textbf{L} — Locks the TLB entry so that it is not considered during replacement candidate selection.
  \item \textbf{U} — Marks the mapped page as uncached
\end{smallGapItemize}

\end{iPageDescEntry}

%--------------------------------------------------------------------------------
\begin{iPageDescOpEntry}{Example}
\begin{lstlisting}[style=iPageOpStyle]

# insert a translation for 0x1000_FFFF_0000 at 0x1_0000 with a 
# page size of 64Kb, access rights read-write, user level

(8)->itlb 0x1000_FFFF_0000, 0x1_0000, 1, 0x7
\end{lstlisting}
\end{iPageDescOpEntry}

%--------------------------------------------------------------------------------
\begin{iPageDescEntry}{Notes}
None.
\end{iPageDescEntry}

%--------------------------------------------------------------------------------
%
%
%--------------------------------------------------------------------------------
\pagebreak
\section*{LF - load ELF file}
\addcontentsline{toc}{section}{LF - load ELF file}

%--------------------------------------------------------------------------------
\begin{iPageDescEntry}{Syntax}
\texttt{lf "<filePath>" } \\
\end{iPageDescEntry}

%--------------------------------------------------------------------------------
\begin{iPageDescEntry}{Description}

The \texttt{LF} command opens an ELF format file and places the ELF segments at the specified address in memory. 

This command is currently under construction and will be finalized when the actual ELF file segments and sections are defined for Twin-64. Currently, it just loads the ELF segments at the specified physical address.

\end{iPageDescEntry}

%--------------------------------------------------------------------------------
\begin{iPageDescOpEntry}{Example}
\begin{lstlisting}[style=iPageOpStyle]

(5)-> LF "~/elf-file"    # load elf-file

\end{lstlisting}
\end{iPageDescOpEntry}

%--------------------------------------------------------------------------------
\begin{iPageDescEntry}{Notes}
None.
\end{iPageDescEntry}

%--------------------------------------------------------------------------------
%
%
%--------------------------------------------------------------------------------
\pagebreak
\section*{MA - modify physical memory}
\addcontentsline{toc}{section}{MA - modify physical memory}

%--------------------------------------------------------------------------------
\begin{iPageDescEntry}{Syntax}
\texttt{MA } \\
\end{iPageDescEntry}

%--------------------------------------------------------------------------------
\begin{iPageDescEntry}{Description}

The \texttt{MA} command ...

\end{iPageDescEntry}

%--------------------------------------------------------------------------------
\begin{iPageDescOpEntry}{Example}
\begin{lstlisting}[style=iPageOpStyle]

... 

\end{lstlisting}
\end{iPageDescOpEntry}

%--------------------------------------------------------------------------------
\begin{iPageDescEntry}{Notes}
None.
\end{iPageDescEntry}

%--------------------------------------------------------------------------------
%
%
%--------------------------------------------------------------------------------
\pagebreak
\section*{MR - modify register}
\addcontentsline{toc}{section}{MR - modify register}

%--------------------------------------------------------------------------------
\begin{iPageDescEntry}{Syntax}
\texttt{MR } \\
\end{iPageDescEntry}

%--------------------------------------------------------------------------------
\begin{iPageDescEntry}{Description}

The \texttt{MR} command ...

\end{iPageDescEntry}

%--------------------------------------------------------------------------------
\begin{iPageDescOpEntry}{Example}
\begin{lstlisting}[style=iPageOpStyle]

... 

\end{lstlisting}
\end{iPageDescOpEntry}

%--------------------------------------------------------------------------------
\begin{iPageDescEntry}{Notes}
None.
\end{iPageDescEntry}

%--------------------------------------------------------------------------------
%
%
%--------------------------------------------------------------------------------
\pagebreak
\section*{NM - create a new T64 module}
\addcontentsline{toc}{section}{NM - create a new T64 module}

%--------------------------------------------------------------------------------
\begin{iPageDescEntry}{Syntax}
\texttt{NM <mType> , <modNum>  [ \{ , <key>=<val> \}* ] } \\
\end{iPageDescEntry}

%--------------------------------------------------------------------------------
\begin{iPageDescEntry}{Description}

A Twin-64 system is a set of modules with at least one processor, one memory and one I7O module. The \texttt{NM} command adds a module to the Twin-64 simulator module table. The command is intended to be used by a configuration file. A configuration file will simply contains a list of commands to setup a system. 

The \texttt{module type} parameter defines the module to be added. There are three types, \texttt{PROC}, \texttt{MEM} and \texttt{IO}. The module number parameter assigns the module number. The Twin-64 system supports up to 64 modules at the system level. 

The remaining optional parameters pass configuration values for the module. They all are of the form "\texttt{name=val}". 

\begin{flushleft}
\begin{tabularx}{\linewidth}{|l|l|X|}
\hline
\textbf{Keyword}  & \textbf{Value} & \textbf{Usage} \\ \hline
\texttt{SPA\_ADR} & \texttt{32-bit address} & Soft physical address start. The address must be aligned to the length parameter. \\ \hline
\texttt{SPA\_LEN} & \texttt{32-bit size} & Soft physical address length. The size is the pageSize multiplied by a power of 2.\\ \hline
\texttt{TLB} & \texttt{TLB\_FA\_xxS} & \ TLB set size. The defined values are 16, 32, 64 and 128. \\ \hline
\texttt{MEM} & \texttt{RAM | ROM} & Memory access type. There are read/write and read only memories. \\ \hline
\end{tabularx}
\end{flushleft}

\end{iPageDescEntry}

%--------------------------------------------------------------------------------
\begin{iPageDescOpEntry}{Example}
\begin{lstlisting}[style=iPageOpStyle]

# add a processor module - 3 - with defaults.
NM proc, 3

# add a processor module - 4 - with a larger TLB                   
NM proc, 4, tlb=TLB_FA_128S

# add a memory module covering the first 256 pages
NM mem, 7, spa_adr=0x0, spa_len=256*0x1000

\end{lstlisting}
\end{iPageDescOpEntry}

%--------------------------------------------------------------------------------
\begin{iPageDescEntry}{Notes}
None.
\end{iPageDescEntry}

%--------------------------------------------------------------------------------
%
%
%--------------------------------------------------------------------------------
\pagebreak
\section*{PTLB — Purge TLB entry}
\addcontentsline{toc}{section}{PDTLB — Purge TLB entry}

%--------------------------------------------------------------------------------
\begin{iPageDescEntry}{Syntax}
\texttt{ptlb [ <vAdr> ]}
\end{iPageDescEntry}

%--------------------------------------------------------------------------------
\begin{iPageDescEntry}{Description}

The \texttt{PDTLB} command invalidates entries in the unified Translation Lookaside Buffer (TLB). If a virtual address \texttt{vAdr} is specified, only the TLB entry matching that virtual address range is purged. If no address is supplied, all data TLB entries are invalidated.

\end{iPageDescEntry}

%--------------------------------------------------------------------------------
\begin{iPageDescOpEntry}{Example}
\begin{lstlisting}[style=iPageOpStyle]

(3)-># purge the entire instruction TLB
(3)->ptlb

(4)-># purge the instruction TLB entry covering virtual address 0x1000_0000
(4)->ptlb 0x80_1000_0000

\end{lstlisting}
\end{iPageDescOpEntry}

%--------------------------------------------------------------------------------
\begin{iPageDescEntry}{Notes}
If the simulator environment consists of serval processors, all TLBs will invalidate the specified virtual address range.
\end{iPageDescEntry}

%--------------------------------------------------------------------------------
%
%
%--------------------------------------------------------------------------------
\pagebreak

\unnumberedSection{REDO - Redo command from history stack}

%--------------------------------------------------------------------------------
\begin{iPageDescEntry}{Syntax}
\texttt{REDO [ <cmdId> ]}
\end{iPageDescEntry}

%--------------------------------------------------------------------------------
\begin{iPageDescEntry}{Description}

The \texttt{REDO} command recalls a previously executed command from the history stack using its command Id and opens it for editing in the interactive command line.

Editing is performed directly in the command line buffer. Characters can be inserted or removed at the cursor position. Pressing \texttt{Enter} submits the modified command for execution.

The \texttt{REDO} command is only available in interactive console sessions (TTY). It is not available when input is read from a script file or redirected from a non-interactive source.

\end{iPageDescEntry}

%--------------------------------------------------------------------------------
\begin{iPageDescOpEntry}{Example}
\begin{lstlisting}[style=iPageOpStyle]

(11) ->hist
 [0]: rmod all
 [1]: nmod tlb,  0, TLB_FA_16S
 [2]: nmod proc, 4
 [3]: nmod mem,  1, ROM, SPA_ADR=0xff_0000_0000, SPA_LEN=0x4000
 [4]: nmod mem,  2, RAM, SPA_ADR=0x0000_0000, SPA_LEN=0x40_0000
 [5]: reset all
 [6]: env halt_on_traps, true
 [7]: loadelf "~/Loop-Test.out"
 (11) ->redo 2
 nmod proc, 4


\end{lstlisting}
\end{iPageDescOpEntry}

%--------------------------------------------------------------------------------
\begin{iPageDescEntry}{Notes}

The command depends on interactive terminal input and is therefore disabled in non-interactive execution modes such as script files or piped input.

\end{iPageDescEntry}
%--------------------------------------------------------------------------------
%
%
%--------------------------------------------------------------------------------
\pagebreak
\section*{RM - remove a T64 module}
\addcontentsline{toc}{section}{RM - remove a T64 module}

%--------------------------------------------------------------------------------
\begin{iPageDescEntry}{Syntax}
\texttt{RM } \\
\end{iPageDescEntry}

%--------------------------------------------------------------------------------
\begin{iPageDescEntry}{Description}

The \texttt{RM} command removes a T64 module from the system. If there are any windows associated with the module, they will be removed as well.

\end{iPageDescEntry}

%--------------------------------------------------------------------------------
\begin{iPageDescOpEntry}{Example}

The following example removes the T64 module number 3.

\begin{lstlisting}[style=iPageOpStyle]
RM 3
\end{lstlisting}
\end{iPageDescOpEntry}

%--------------------------------------------------------------------------------
\begin{iPageDescEntry}{Notes}
None.
\end{iPageDescEntry}

%--------------------------------------------------------------------------------
%
%
%--------------------------------------------------------------------------------
\pagebreak
\section*{RESET - Reset T64 modules}
\addcontentsline{toc}{section}{RESET - Reset T64 modules}

%--------------------------------------------------------------------------------
\begin{iPageDescEntry}{Syntax}
\texttt{RESET <modNum> | ALL } \\
\end{iPageDescEntry}

%--------------------------------------------------------------------------------
\begin{iPageDescEntry}{Description}

The \texttt{RESET} command resets a module or the entire system. The module will execute its defined reset tasks. The result is module specific. A memory module will clear its memory range, an I/O module will reset its internal state, and a processor will clear its internal registers and start at the architected instruction address.

Like all commands dealing with module management, this command should be taken with care. A common use could be to have this command as the last command in a configuration file to start the Twin-64 system.

\end{iPageDescEntry}

%--------------------------------------------------------------------------------
\begin{iPageDescOpEntry}{Example}
\begin{lstlisting}[style=iPageOpStyle]

RESET 3               # reset module number three.

\end{lstlisting}
\end{iPageDescOpEntry}

%--------------------------------------------------------------------------------
\begin{iPageDescEntry}{Notes}
None.
\end{iPageDescEntry}

%--------------------------------------------------------------------------------
%
%
%--------------------------------------------------------------------------------
\pagebreak
\section*{RUN — Run the simulation}
\addcontentsline{toc}{section}{RUN — Run the simulation}

%--------------------------------------------------------------------------------
\begin{iPageDescEntry}{Syntax}
\texttt{RUN} \\
\end{iPageDescEntry}

%--------------------------------------------------------------------------------
\begin{iPageDescEntry}{Description}

The \texttt{RUN} command starts continuous execution of the simulation. It is
functionally equivalent to repeatedly stepping the simulation without an
explicit step count.

Once invoked, control is transferred to the TWIN-64 system, and all configured
and enabled processors execute instructions continuously. Execution proceeds
until a condition occurs that returns control to the simulator.

Control is returned to the simulator when a trap, breakpoint exception, trace
exception, machine check, or other architectural exception is raised. At that
point, execution halts and the simulator command interface becomes active again,
allowing inspection or modification of system state.

Unlike the \texttt{STEP} command, \texttt{RUN} does not impose a limit on the
number of executed instructions and continues indefinitely until an exceptional
event occurs.

\end{iPageDescEntry}

%--------------------------------------------------------------------------------
\begin{iPageDescOpEntry}{Example}
\begin{lstlisting}[style=iPageOpStyle]

(80)->run          # start continuous simulation execution

\end{lstlisting}
\end{iPageDescOpEntry}

%--------------------------------------------------------------------------------
\begin{iPageDescEntry}{Notes}

The \texttt{RUN} command respects breakpoint and trace settings. If breakpoints or
single-step tracing are enabled, execution may stop shortly after \texttt{RUN}
is issued.

\end{iPageDescEntry}

%--------------------------------------------------------------------------------
%
%
%--------------------------------------------------------------------------------
\pagebreak
\unnumberedSection{STEP - Step a T64 module}

%--------------------------------------------------------------------------------
\begin{iPageDescEntry}{Syntax}
\texttt{STEP [ <steps> [ , <modNum> ]} \\
\end{iPageDescEntry}

%--------------------------------------------------------------------------------
\begin{iPageDescEntry}{Description}

The \texttt{STEP} command advances a halted T64 module by executing one or more 
instructions on a processor module or step units on an I/O module. There can be 
one or more processors and I/O modules configured. When the step command is 
invoked without an argument, the simulator executes exactly one instruction on
all modules and return s control to the simulator. Likewise, a step count greater
than one will execute that amount of steps. 

If the module number is provided and the module is halted, only the module specified 
will execute the entered steps.

In either case, after the requested number of steps have completed, execution 
halts and control returns to the simulator command interface, allowing inspection 
of state, registers, memory, or windows. After inspection the step command or 
the run command resumes execution of the processor modules.

\end{iPageDescEntry}

%--------------------------------------------------------------------------------
\begin{iPageDescOpEntry}{Example}
\begin{lstlisting}[style=iPageOpStyle]

(5)->s        # execute a single instruction on all modules.
(6)->

(6)->s 10     # execute ten instructions in all modules.
(7)->

(8)->s 10, 5  # execute ten instructions on module 5
(9)->

\end{lstlisting}
\end{iPageDescOpEntry}

%--------------------------------------------------------------------------------
\begin{iPageDescEntry}{Notes}

The \texttt{STEP} command respects breakpoint and trace settings. A breakpoint 
or trace exception encountered during stepping may cause execution to stop before
the requested number of steps has completed.

\end{iPageDescEntry}

%--------------------------------------------------------------------------------
%
%
%--------------------------------------------------------------------------------
\pagebreak 
\unnumberedSection{W - Display an expression value} 

%------------------------------------------------------------------------------
\begin{iPageDescEntry}{Syntax} 
\texttt{W <expr>} \\ 
\end{iPageDescEntry}

%--------------------------------------------------------------------------------
\begin{iPageDescEntry}{Description}

The \texttt{W} command evaluates a simulator expression and displays its
current value. Expressions can include:

\begin{smallGapItemize}
  \item Registers and control/status registers
  \item Memory locations (by virtual or physical address)
  \item Arithmetic operations, logical operations, and constants
  \item Predefined functions provided by the simulator
\end{smallGapItemize}

The command is primarily used for inspecting system state during simulation
or debugging. The evaluated value is printed in the current radix or display
format of the active window. If no expression is provided, the command may
display a prompt or an error, depending on the simulator configuration.

\end{iPageDescEntry}

%--------------------------------------------------------------------------------
\begin{iPageDescOpEntry}{Example}
\begin{lstlisting}[style=iPageOpStyle]

# display the value of general-purpose register R1
(12)-> w R1
0x0000_0000_1234_5678

(13)-> w 0x100, dec
 256

# display the value at virtual memory address 0x1000
(14)-> w [0x1000]
0xDEAD_BEEF

# evaluate an arithmetic expression
(15)-> w R1 + R2 * 4
0x0000_0000_2468_ACF0

\end{lstlisting}
\end{iPageDescOpEntry}

%--------------------------------------------------------------------------------
\begin{iPageDescEntry}{Notes}

None.

\end{iPageDescEntry}

%--------------------------------------------------------------------------------
%
%
%--------------------------------------------------------------------------------
\pagebreak
\section*{XF - execute commands from a file}
\addcontentsline{toc}{section}{XF - execute commands from a file}

%--------------------------------------------------------------------------------
\begin{iPageDescEntry}{Syntax}
\texttt{XF } \\
\end{iPageDescEntry}

%--------------------------------------------------------------------------------
\begin{iPageDescEntry}{Description}

The \texttt{XF} command ...

\end{iPageDescEntry}

%--------------------------------------------------------------------------------
\begin{iPageDescOpEntry}{Example}
\begin{lstlisting}[style=iPageOpStyle]

... 

\end{lstlisting}
\end{iPageDescOpEntry}

%--------------------------------------------------------------------------------
\begin{iPageDescEntry}{Notes}
None.
\end{iPageDescEntry}

